\section{The performant configuration}\label{sec:config}

Everything above is the construction. This appendix is what to actually set,
and why. Values are read from the repository as committed; measurements are
cited where they exist and labelled where they do not.

\subsection{Parameters}

\begin{center}\small
\begin{tabular}{@{}p{0.16\textwidth}p{0.15\textwidth}p{0.30\textwidth}@{}}
\toprule
parameter & value & why\\
\midrule
commitment field & $\mathrm{GF}(2^{128})$, GHASH & fixed by the packing width;
  $\mathrm{GF}(2^{127})$ measures $0$--$12\%$ \emph{slower} on this repo's hot
  patterns\\
\code{LOG\_PACKING} & $7$, \textbf{not tunable} & the ring-switch recombination
  is $128\times128$ over $\Ftwo$; the packing count must be a power of two
  matching the field dimension\\
Ligerito \code{initial\_k} & $\mathbf 4$ & the knee at every rate: $k{=}6\to4$
  is $399.9\to237.8$\,KiB ($-40\%$) for $+3\%$ prove time. $k{=}3$ adds a
  ladder level and backfires\\
Ligerito base rate & $1/8$ & $+5\%$ prove for $-35\%$ proof against rate $1/2$;
  past $1/16$ the Ligerito blob is smaller than the forest side anyway\\
evaluation modulus & $q=2^{100}-15$, fixed & \cref{warn:fixedq}; keeps the
  chunk count at $1$ for $t\leq26$\\
\code{word\_bits} $W$ & $\mathbf 1$ & the chunk width is $127-t-W$; at $W{=}32$
  it drops below $100$ for every $t$, forcing a second chunk and doubling the
  forest\\
\bottomrule
\end{tabular}
\end{center}

\paragraph{Layout and claim routing.}
\begin{center}\small
\begin{tabular}{@{}p{0.16\textwidth}p{0.15\textwidth}p{0.30\textwidth}@{}}
\toprule
parameter & value & why\\
\midrule
committed columns & $8$ padded, taps layout & the bit-diagonal layout cannot
  express rotations at all --- see below\\
word group $g$, axis & $g{=}5$, in the \textbf{low} $5$ bits of the clear axis
  & hard assertion $g\leq s$; $\ROTR$/$\SHR$ are permutations of clear rows\\
clear width $s$ & $s=\lfloor n/2\rfloor$; \textbf{needs $n\geq22$} & sets the
  stackable-claim ceiling $R_{\max}=2^{\,s-g}-1$ (\cref{warn:envelope}); at
  $n{=}21$ that is $31<48$ and the schedule is \emph{silently clipped}\\
forest depth $t'$ & cap at $16$--$17$ and hold constant as $N$ grows & the cache
  cliff; do \emph{not} use the shipped $t\approx0.6n$\\
$x$-fold extra $\delta$ & $3$--$4$ on XOR bodies, $5$ on the plain layer &
  measured: $86.4\to71.2$\,ms and $284\to163$\,KB at $n{=}24$\\
claim routing & multiweight collapse for linear, composed collapse for the
  schedule, blocked tap-claims for XOR shapes & \textbf{never} the batched
  virtual-XOR path\\
\code{F2Z\_QUAD} & \textbf{off} & $-12/-20\%$ at $n{=}22/24$ but $+6/+5\%$ at
  $n{=}26/28$\\
\bottomrule
\end{tabular}
\end{center}

\subsection{The one thing to get right}

The F2Z paper's §4.1 opens with the author's own note that \emph{``in current
implementation, this performs worse than not using any virtualization, which
should not be the case.''} That anomaly has a diagnosis, and it is not a
constant factor.

There are two virtualization surfaces in the repository. The one the note's
measurement uses is \textbf{diagonal in the bit-position index}: source offset
and destination share the same $j$, so it XORs across \emph{columns} at a fixed
bit position \src{extract\_virtual\_xor\_rows --- src/pcs.rs}. A rotation
permutes $j$. \textbf{That path cannot express $\Sigma_0$, $\Sigma_1$, $\sigma_0$
or $\sigma_1$ at all}, so the measurement was taken on a surface that cannot run
SHA-256. Worse, it runs the main committed claim's forests \emph{and} a batched
extra forest on top --- strictly more grand-product work than not virtualizing.

The surface that does work is the tap machinery \src{src/taps.rs}, which crosses
lanes and requires the word field inside the clear axis --- the layout pin of
\cref{rem:layout}. Routing claims through the collapse and shared-point families
instead of the batched path measures, at $n{=}26$ with $15$ claims:

\begin{center}\small
\begin{tabular}{@{}lrrr@{}}
\toprule
& batched virtual-XOR & collapse / RLC & gain\\
\midrule
prover & $967$\,ms & $\mathbf{431}$\,ms & $2.2\times$\\
proof & $1193$\,KB & $\mathbf{272}$\,KB & $4.4\times$\\
verify & $33.1$\,ms & $\mathbf{11.0}$\,ms & $3.0\times$\\
\bottomrule
\end{tabular}
\end{center}

\noindent
It also removes a hard memory wall: the batched path pads to $17.2$\,GB at
$n{=}28$, which is unrunnable.

\paragraph{An independent reproduction.} A separate run, on a different claim
set --- the harness's own \code{-{}-taps vx} instance ($6$ claims) against
\code{-{}-taps sched} at $48$ --- reproduces the prover figure and shows it is
the large-$n$ limit rather than a fixed constant:

\begin{center}\small
\begin{tabular}{@{}rrrrrrr@{}}
\toprule
& \multicolumn{2}{c}{prove (ms)} & & \multicolumn{2}{c}{proof (KiB)} &\\
\cmidrule(lr){2-3}\cmidrule(lr){5-6}
$n$ & batched & composed & gain & batched & composed & gain\\
\midrule
$22$ & $192.7$  & $61.5$  & $3.13\times$ & $371.3$ & $206.8$ & $1.80\times$\\
$24$ & $430.4$  & $166.9$ & $2.58\times$ & $579.7$ & $283.4$ & $2.05\times$\\
$26$ & $1050$   & $450.0$ & $2.34\times$ & $990.1$ & $433.6$ & $2.28\times$\\
\bottomrule
\end{tabular}
\end{center}

\noindent
The prover gain \emph{decreases} toward $\approx2.2\times$ as $n$ grows while the
proof gain \emph{increases}; quoting either without its $n$ is meaningless. Per
claim the gap is far wider --- at $n{=}22$ the batched path costs
$61.9$\,KiB/claim against the composed path's $4.31$.

\begin{center}\small
\begin{tabular}{@{}lrrrr@{}}
\toprule
mode ($n{=}22$) & $k$ & prove (ms) & proof (KiB) & KiB/claim\\
\midrule
\code{family} clustered stream & $6$  & $410.7$ & $366.9$ & $61.16$\\
\code{vx} batched taps         & $6$  & $192.7$ & $371.3$ & $61.88$\\
\code{collapse} single-tap     & $13$ & $128.4$ & $268.4$ & $20.65$\\
\code{rotxor} op-of-XOR-set    & $8$  & $111.0$ & $268.3$ & $33.53$\\
\code{sched} composed          & $48$ & $\mathbf{61.5}$ & $\mathbf{206.8}$ & $\mathbf{4.31}$\\
\bottomrule
\end{tabular}
\end{center}

\begin{remark}[Reproducing these numbers]\label{rem:repro}
The repository does not build as committed: \code{Cargo.toml} pins
\code{flock-core} to an absolute path outside the repo, the default toolchain
rejects flock's unstable \code{isolate\_lowest\_one}, and current flock has
removed the BaseFold backend and added a \code{merkle\_hash} field. With those
reconciled, every figure above comes from
\code{RUSTFLAGS="-C target-cpu=native" cargo +1.97.1 build -{}-release
-{}-features unchecked}, medians of three, each repetition verified in-process.
The \code{-{}-features unchecked} and \code{target-cpu=native} flags are
load-bearing on \code{aarch64} and the binary says so itself.
\end{remark}

\begin{warning}[The machinery is landed; the routing is not]
The collapse and RLC paths exist, are tested and are transcript-pinned --- $112$
library tests pass, including round-trip and tamper-rejection tests for each of
the tap, collapse, composed and multiweight paths. What does not exist is a claim
planner that groups SHA-256's virtual claims by shared point, column group and
canonical tap source, and the SHA harness that would use it. That is the highest
gain-per-effort item available, and the virtual-XOR driver it would replace
currently has no test coverage.
\end{warning}

\subsection{The composed collapse, measured}\label{sec:measured}

The claim above --- that shift-invariant workloads cost \emph{two} inner bodies
rather than one per claim --- is the load-bearing one for SHA-256, because the
message schedule is exactly such a workload:
\[
  W[t]=\sigma_1(W[t-2])+W[t-7]+\sigma_0(W[t-15])+W[t-16],
  \qquad t=16,\dots,63 .
\]
All $48$ rounds are one read pattern displaced by one word each time, so they
share a single canonical source and differ only in the outer offset $\off^t$.

It is worth being precise about what collapses. \textbf{The witness does not.}
There is one committed bit vector throughout, and all $48$ rounds are statements
about different slices of it. What the $\gamma$-combination merges is the
\emph{claims}: $\ipq{\ww_i}{\text{data}}=c_i$ for $i=1,\dots,k$ becomes the single
$\ipq{\sum_i\gamma_i\ww_i}{\text{data}}=\sum_i\gamma_ic_i$. That is only useful
because the combined weight vector \emph{keeps the tensor shape} the commitment
can answer --- which it does precisely because the $48$ are one shifted pattern:
\[
  \sum_i \gamma_i\,\ww_i
  \;=\; \ww_{\mathrm{row}}^{(0)}\otimes E_0 \;+\; \ww_{\mathrm{row}}^{(1)}\otimes E_1,
  \qquad E_\beta=\sum_i \gamma_i\,e_i^{(\beta)} \bmod q .
\]
Two terms, because a displaced read either stays inside its $32$-bit word
($\beta{=}0$) or spills into the next ($\beta{=}1$). There is no third case.
In the implementation this is literal: the planner de-duplicates claims to
\emph{(source, branch)} pairs and therefore emits exactly two inner claims at any
$k$ \src{tap\_composed\_plan --- src/ligerito\_flock.rs}, and the prover's inner
list is the \emph{plan}, not the claims.

Sweeping $k$ to the envelope limit at $n{=}26$ confirms it
(\code{-{}-taps sched -{}-taps-rounds $k$}, median of $3$, every rep verified):

\begin{center}\small
\begin{tabular}{@{}rrrrrr@{}}
\toprule
$k$ & prove (ms) & verify (ms) & proof (KiB) & KiB/claim & verify/prove\\
\midrule
$1$   & $219.7$ & $16.8$  & $299.7$ & $299.70$ & $0.08$\\
$2$   & $371.1$ & $31.5$  & $434.1$ & $217.04$ & $0.08$\\
$16$  & $386.2$ & $49.4$  & $433.7$ & $27.10$  & $0.13$\\
$48$  & $404.3$ & $89.8$  & $433.6$ & $9.03$   & $0.22$\\
$128$ & $384.6$ & $191.2$ & $433.9$ & $3.39$   & $0.50$\\
$255$ & $420.7$ & $348.8$ & $433.7$ & $1.70$   & $\mathbf{0.83}$\\
\bottomrule
\end{tabular}
\end{center}

\noindent
Proof size steps \emph{once} --- $299.7\to434.1$\,KiB between $k{=}1$ and
$k{=}2$, which is the second branch appearing --- and then does not move through
$k{=}255$. That is $127\times$ the claims for $0.5$\,KiB. Prove time is flat.
Against the counterfactual of proving the $48$ schedule rounds one at a time,
$48\times219.7$\,ms $=10.5$\,s and $48\times299.7$\,KiB $=14.0$\,MiB become
$404$\,ms and $434$\,KiB: $\mathbf{26\times}$ prove and $\mathbf{33\times}$ proof.

\begin{warning}[The verifier is the term that grows, not the prover]
\label{warn:verifyk}
The collapse is free on the prover and free on proof size, but the verifier must
absorb every claim's statement and form $E_\beta=\sum_i\gamma_ie_i^{(\beta)}$,
which is $\Theta(k)$ --- $O(k\cdot2^{\,s-\delta})$ in the implementation
\src{src/ligerito\_flock.rs}. The measured ratio climbs $0.08\to0.83$ over
$k{=}1\dots255$, so \textbf{verifier cost overtakes prover cost near
$k\approx300$}. At SHA-256's $48$ schedule rounds the ratio is $0.22$ and there
is ample room; a planner that stacks many claim \emph{families} into one group
must watch this number, because it is the only cost in the collapse that is not
amortised. This inverts the usual concern: what limits stacking is the verifier.
\end{warning}

\begin{warning}[The stackable-claim envelope, and a silent clip]
\label{warn:envelope}
Word offsets must satisfy $g\leq s$ and $\off<2^{\,s-g}$
\src{src/taps.rs}, and the harness derives $s=\lfloor n/2\rfloor$. The number of
schedule-shaped claims that can be stacked is therefore
\[
  R_{\max}(n,g)\;=\;2^{\,\lfloor n/2\rfloor-g}-1,
  \qquad\text{independent of }\delta ,
\]
measured as $63$, $255$, $1023$ at $n{=}22,26,30$ with $g{=}5$, and confirmed
against the assertions. SHA-256 needs $R\geq48$, so \textbf{$n\geq22$}; at
$n{=}21$ the ceiling is $31$ and the harness \emph{silently clips} the requested
$48$ rather than failing. Any A/B taken at $n{=}21$ or below is therefore not
measuring the full schedule. Note also that the reported KiB/claim figure falls
as $1/k$ by construction --- it is a divisor, not a measurement --- so it should
never be quoted without the flat proof size beside it.
\end{warning}

\subsection{Where the time goes}

The forest's share of prover time depends strongly on shape. \textbf{Measured}
on the shipped row-heavy split it is $82\%$ at $n{=}26$, $92\%$ at $n{=}28$ and
$95\%$ at $n{=}30$; but only $53$--$65\%$ at $n{=}22$ on the \emph{column-heavy}
geometry that SHA-256's layout pin ($s\geq11$) actually forces, with the
recursive-Ligerito tail taking $17$--$35\%$. So $|\hh|$ is an excellent proxy for
prover time at production size and a mediocre one at the sizes where the tap
machinery was tuned --- which means small-$n$ A/Bs will systematically
\emph{understate} the value of $|\hh|$ levers. Either way, virtualization shrinks
what is \emph{committed}, not what is \emph{folded}:

\begin{center}\small
\begin{tabular}{@{}p{0.24\textwidth}p{0.28\textwidth}p{0.24\textwidth}@{}}
\toprule
stage & scales with & virtualization helps?\\
\midrule
witness generation & $N\cdot(|\ff|+|\hh|)$ & no\\
commit, encode, Merkle & $N\cdot|\ff|$ & \textbf{yes}\\
PIOP evaluation-shaping sumcheck & $N\cdot|\hh|$ & \textbf{no} --- one required pass\\
exponent fold $+$ GKR forest & $N\cdot|\hh|$ & \textbf{no} --- the bottleneck\\
GKR-boundary pre-sumcheck & $N\cdot|\hh|$ & \textbf{no}\\
ring-switch & $N\cdot|\ff|$ & yes\\
Ligerito open & $N\cdot|\ff|$ & yes\\
public row collapse (\cref{step:collapse}) & one sparse matrix pass; public-I/O reading is $\Theta(N)$ & n/a\\
\bottomrule
\end{tabular}
\end{center}

So the committed-size win lands on the smaller half of the budget. The
open question that would move the larger half is the paper's own: \emph{``in
reality many of these bits are linear combinations in $\Ftwo$ of witness
bits''} --- exploiting that structure \emph{inside} the grand product. Nobody has
done it.

\subsection{Not worth doing}

\begin{itemize}[leftmargin=1.4em,itemsep=2pt,topsep=4pt]
  \item \textbf{$\mathrm{GF}(2^{127})$.} Attractive on paper --- an irreducible
        trinomial gives a PMULL-free reduction, and $2^{127}-1$ is a Mersenne
        prime so every $\alpha\notin\{0,1\}$ generates and the primitivity test
        disappears. But it is architecturally excluded (the ring-switch needs a
        power-of-two field dimension) and measures slower on Apple silicon,
        where PMULL throughput makes GHASH's all-multiply pipeline the better
        instruction mix.
  \item \textbf{A rank-$r$ decomposition of the original reshaped $\vv$.}
        It repairs the invalid rank-one shortcut, but repeats the integer-fold,
        forest, and GKR weight families $r$ times.  The required
        evaluation-shaping sumcheck of \cref{step:eqshape} costs one
        $\Theta(N|\hh|)$ pass and produces rank-one equality weights for every
        axis split.
  \item \textbf{Raising the forest depth with $n$.} The shipped
        $t\approx0.6n$ walks off a cache cliff; cap $t'$ and let the tree count
        absorb growth.
\end{itemize}

\subsection{Open items}

\begin{enumerate}[leftmargin=1.5em,itemsep=2pt,topsep=4pt]
  \item \textbf{Per-column typing} (\cref{rem:typing}) --- without it the
        committed-size gain is zero. The single biggest risk.
  \item \textbf{The fixed-prime side condition} (\cref{warn:fixedq}) --- not a
        lemma and not a gap: the error is zero unconditionally, because the
        residual is bounded by a public constant. What must be recorded is the
        parameter check $q>2\max_{\mathbf b}\sum_j|A_{\mathbf b,j}|\approx2^{38}$,
        to be re-derived if $A$ changes. \emph{Downgraded from a soundness gap.}
  \item \textbf{Implement or fuse the evaluation-shaping sumcheck}
        (\cref{step:eqshape,warn:rank1}).
        The protocol now resolves the mathematical gap before the opening:
        it sumchecks the arbitrary Step 3 coefficient vector into
        $c\,\eq(\cdot,r_h)$.
        \begin{itemize}[leftmargin=1.2em,itemsep=1pt,topsep=2pt]
          \item Directly reshaping the Step 3 coefficient matrix $V$ remains
                invalid: its measured rank along the shipped split is
                $1,2,\mathbf 4,4,3,2$ for $s-g=0,\dots,5$ against a
                schedule-shaped $A_0$, with $s=\lfloor n/2\rfloor$ shipped.
          \item The implementation must add this sumcheck, or fuse it with the
                sparse row collapse as a Spartan-style matrix sumcheck, and
                include the extra $\Theta(N|\hh|)$ pass in benchmarks.
        \end{itemize}
        The previous rank-one assertion was a \emph{completeness} failure; it
        must not survive as an implementation shortcut.
  \item \textbf{Prove the clear-dependent tap closure equal to the full
        adjoint} (\cref{step:transpose,warn:mblockdiag}).
        \Cref{step:transpose} states the always-correct
        $M^{\!\top}q_H$.  The old shortcut
        $M=\mathrm{Id}_{\text{col}}\otimes M_0$ is false for SHA-256, so the tap
        implementation still needs a protocol-matching equivalence proof.
  \item \textbf{$\ellcol=2^{s}$, and the $\Theta(N)$ proof-size claim was wrong}
        --- now resolved against the code, recorded here because it propagated.
        Three verifier sites hard-reject unless the sent-integer count equals
        $2^{s}$ exactly, and $N$ is not a parameter of \code{IntEvalParams} at
        all. Since every layout sets $s$ to a fixed fraction of $n$, the block is
        $\Theta(N^{2/5})$--$\Theta(N^{1/2})$. \Cref{warn:lanes} records the
        corrected axis interpretation; the old $\ellcol=N$ argument has been
        removed.
  \item \textbf{The missing theorem is the \emph{virtualized} one --- and that is
        the one this note uses.} Earlier drafts of this note claimed the paper
        had no proved soundness for the core IOP at all. \textbf{That was read
        off a stale checkout and is false}, and the correction is a narrowing:
        \begin{itemize}[leftmargin=1.2em,itemsep=1pt,topsep=2pt]
          \item the \emph{core}-IOP theorem is stated with explicit hypotheses
                (including $\ord(\gen)>(\ellrow{+}1)(q_{\max}{-}1)$) and carries a
                complete proof --- knowledge-state function, round-by-round
                extractor, its own reduction lemma. Only \emph{completeness} is
                flagged missing on it;
          \item the PIMSAT construction's theorem is stated with errors and
                proved except for one paragraph;
          \item the \emph{virtualized} variant is still an empty
                statement--proof pair. There are no \code{WIP} markers in any
                proof in the paper; the two that exist are in the introduction's
                benchmark prose.
        \end{itemize}
        So the hole is real but sits exactly where this note needs it, since
        everything here runs through the virtualized relation. \textbf{Still the
        largest hole on this list} --- for a narrower and checkable reason.
  \item \textbf{Bit-ness as a soundness precondition} (\cref{warn:bitness}) ---
        the paper flags this itself (\emph{``we started with $\ff$ containing
        only bits, so we need to add a wrapper around it''}, marked \code{todo})
        but does not close it. The argument in \cref{warn:bitness} is ours. Note
        that the wrapper is required for a \emph{module} witness and not for the
        bit witness used here.
  \item \textbf{Public I/O still costs $\Theta(N)$.}  The h-only normalization
        removes a separate public-fold protocol step by substituting $\xx$ into
        $A_{\xx}$ and carrying one constant coordinate.  It does not make $N$
        independent public messages and digests free: the verifier must still
        read them when constructing $A_{\xx}$.  The sent-integer block is
        $\ellcol=2^s$, not $N$.
  \item \textbf{The homogeneous h-only relation is a local normalization, not
        a relation stated by the paper.}  This note uses
        $\ff^{\circ}=(1,\ff)$,
        $\pitwo(\hh)=M\pitwo(\ff^{\circ})$, and
        $A_{\xx}\hh=0$.  §4.1 gives only a single inner-product claim.  The
        constant coordinate, public substitution, and exact implementation
        index map must therefore be proved equivalent to the SHA-256 relation
        rather than cited as an existing paper construction.
  \item \textbf{The paper's own virtualization target is unreachable.}
        \code{paper/main.tex} states that virtualizing XOR should take the
        witness from $2^{14+k}$ to $2^{12+k}$ for $2^k$ hashes. The census gives
        $|\ff|=6\,816>4\,096$, so it pads to $2^{13}$ --- missing the stated
        target by $1.66\times$ while beating the unvirtualized $2^{14}$ baseline
        by $2.4\times$. Neither document cites the other on this point.
  \item \textbf{Feed-forward.} Zinc\raisebox{0.1ex}{+} omits it, and the local
        \code{paper/main.tex} still reads \emph{``we omit this step for
        brevity''} --- but the current Overleaf render includes it as
        $H^{\mathrm{out}}$, so the local checkout is stale. Adding it costs eight
        rows and eight one-bit carries, entry set unchanged. The
        $206$-row count elsewhere is for the $64$-round loop alone.
\end{enumerate}
